\documentclass{UAmathtalk}

\author{Tianyi Yu}
\address{Université du Québec à Montréal, Canada}
\urladdr{https://tianyi-math.github.io}
\title{A Positive Formula Involving Pipes and Bruhat Order}
\date{Thursday, September 10, 2026}

\renewcommand*{\where}{Massry B008}

\begin{document}

\maketitle

\begin{abstract}
Positivity phenomena are common in algebraic geometry: coefficients coming from complicated expansions often turn out to be nonnegative integers, suggesting that they should count something. In this talk, I will discuss one such problem arising from Schubert calculus. Lam, Lee, and Shimozono proved certain coefficients in their study of affine Grassmannians are positive and asked for a combinatorial explanation. I will present the first such formula, using two combinatorial models: bumpless pipe dreams, which are configurations of pipes on a grid, and the Bruhat order. I will also explain how a correspondence between these models leads to the positivity formula. This is joint work with Jack Chou.
\end{abstract}

\end{document}
